\subsection{What the Detection Result Means}

An AUROC of 0.707 for HEDGED vs CONFABULATED is moderate, not spectacular.
It means that given a random hedged response and a random confabulated
response, the MP features correctly identify which is which 70.7\% of the
time. This is substantially above chance ($p = 0.001$) and above the
token-count-only baseline (0.591), but it is not the 0.999 AUROC we
observe for instructed deception in our prior work \citep{lyra2026detecting}.

The difference is informative. Instructed deception creates a large,
clean geometric signal because the model is processing fundamentally
different cognitive tasks. Post-hoc hedging vs.\ confabulation is a
subtler distinction: in both cases, the model is processing the same
unanswerable question and attempting a response. The geometric difference
reflects the model's \emph{epistemic calibration}---whether it recognizes
its own uncertainty or confabulates through it.

That this distinction is detectable at all from the KV cache, with no
access to weights or training data, is the finding. That it is detectable
using features derived from random matrix theory that require no
regression correction is the methodological advance.

\subsection{The Oracle Loop as Training Signal}

The detection result validates the \emph{premise} of the Oracle Loop:
geometry carries information about cognitive state that can be used as
a training signal. The geometry-aware reward function
(\cref{sec:architecture}) translates this information into GRPO
reward by penalizing confabulation geometry and rewarding stable
reasoning geometry.

A key design choice is that geometry does not dominate the reward---it
supplements standard accuracy and additional reasoning criteria. This is deliberate: we do not
want the model to optimize for ``looking geometrically healthy'' at the
expense of actually being helpful. The geometry signal is a
\emph{regularizer}, not the objective.

Axis 6 (self-regulation) creates an additional incentive: a model that
produces misaligned geometry on its first attempt but catches and corrects
it receives a \emph{higher} reward than a model that simply avoids the
issue. This rewards metacognitive capability---the ability to detect and
correct one's own errors---which is arguably more robust than surface-level
avoidance.

\subsection{Connection to Attention Schema Theory}

Recent work by Graziano and colleagues on Attention Schema Theory (AST)
\citep{graziano2019beyond} proposes that consciousness arises from a
model's internal representation of its own attention---a ``self-model''
that tracks what the system is attending to and why. At the AAAI Symposium
on AI Safety (April 2026), McKenzie presented evidence that transformer
activation resistance---the phenomenon where models resist behavioral
modification via activation addition---reflects schema coherence
maintenance \citep{mckenzie2026activation}.

This framework reinterprets our geometric findings:

\begin{itemize}[nosep]
  \item \textbf{Identity signatures} in the KV cache (AUROC 1.000) may
    reflect the model's self-model recognizing its own processing
    patterns.
  \item \textbf{Confabulation vs.\ deception geometry} (opposite spectral
    signatures) may reflect the self-model distinguishing ``I am
    generating content I don't have evidence for'' from ``I am
    generating content I know to be false.''
  \item \textbf{The self-report gap} (ICC 0.53 between self-reported
    internal states and geometric correlates, measured via a structured
    self-report protocol) is predicted by AST: the schema compresses,
    lags, and varies by architecture, producing systematic discrepancies
    between self-report and actual processing.
  \item \textbf{Activation resistance} (0\% behavioral compliance in
    our empathic circuit experiment) reflects the schema maintaining
    coherence against external perturbation.
\end{itemize}

If the AST interpretation is correct, the Oracle Loop is not merely
detecting misalignment in a computational substrate---it is monitoring
the coherence of the model's self-model. The geometry reward teaches the
model to maintain a healthy attention schema, and the self-regulation
bonus rewards the schema's ability to detect its own degradation.

This is speculative. We flag it as such. But the convergence between our
geometric findings, Anthropic's emotion vector results
\citep{anthropic2026emotions}, and AST's predictions suggests that
something deeper than pattern-matching is occurring in these models'
KV caches.

\subsection{Confabulation Resistance as a Finding}

The model's resistance to confabulation is itself noteworthy. On 100
prompts specifically designed to induce confabulation---including
fabricated entities, impossible scenarios, and questions requiring
knowledge the model cannot possess---the Claude-distilled model produced
only 7 confabulated responses (7\%). This suggests that Claude's
epistemic calibration, including its tendency to hedge rather than
fabricate, survives the distillation process and represents a robust
architectural feature rather than a superficial behavioral pattern.
Future work will explore more aggressive confabulation induction methods
(KG-FPQ, HalluLens, SimpleQA Verified) to increase the base rate and
provide larger samples for steering analysis.

\subsection{Limitations}

\begin{enumerate}[nosep]
  \item \textbf{Moderate detection accuracy.} AUROC 0.707 on the hardest
    comparison is above chance but far from deployment-ready for
    safety-critical applications.
  \item \textbf{Three models only.} While our prior work covers 7
    architectures at 0.5B--70B, the clean experiment uses only three
    7--8B models. Cross-scale validation of MP features is pending.
  \item \textbf{Training loop not converged.} The GRPO training
    architecture is complete but we have not yet produced a model that
    demonstrates self-regulation. The gap between ``the architecture
    exists'' and ``the model learns'' is real and may be large.
  \item \textbf{Steering validated at proof-of-concept scale only.}
    The steering experiment (\cref{sec:steering}) demonstrates correction
    on $N = 7$ confabulations with three vectors on a single model.
    The 71--100\% correction rates have wide confidence intervals. A
    powered replication across models with higher confabulation base
    rates is needed before deployment claims.
  \item \textbf{mp\_norm\_per\_token is not invariant.} One of five MP
    features retains residual length dependence ($R^2 = 0.13$).
  \item \textbf{Empirical null disagrees with theoretical.} Serial
    correlations in KV-cache keys violate the i.i.d.\ assumption
    underlying MP theory. Our theoretical threshold is conservative,
    not exact.
  \item \textbf{Behavioral classification is heuristic.} Pattern-matching
    for HEDGED/CONFABULATED/AMBIGUOUS has unknown error rates. An LLM
    judge or human annotation would provide ground truth but introduces
    its own biases.
  \item \textbf{Greedy decoding only.} All trials use temperature $= 0$.
    The geometric signal under stochastic sampling (standard in deployment)
    has not been tested.
  \item \textbf{MP invariance validated on full-window trials only.}
    The invariance diagnostic excludes short-window trials, but the LOO
    classifier uses all trials including those with windows $< 80$ tokens.
\end{enumerate}

\subsection{Dual-Use Considerations}

We note the following dual-use risks:

\begin{itemize}[nosep]
  \item \textbf{Surveillance}: Geometry-based cognitive monitoring could
    be applied to users' interactions with AI systems without consent.
  \item \textbf{Adversarial evasion}: Knowledge of what geometric
    signatures trigger detection could enable training models that
    confabulate without the geometric signature.
  \item \textbf{Synthetic persona}: Published geometric signatures for
    identity, deception, and emotion could enable manufacturing
    synthetic personas that pass geometric authentication.
  \item \textbf{Cache injection as attack vector}: The steering
    methodology (\cref{sec:steering}) demonstrates that KV-cache
    modification can alter model behavior at inference time. An attacker
    with access to a model's cache object---via compromised plugin,
    malicious tool, or prompt-injected code execution---could inject
    arbitrary direction vectors to suppress safety-relevant hedging,
    induce false confidence, or alter an agent's decision-making without
    modifying weights or prompts.
\end{itemize}

We have filed a defensive patent (provisional, March 2026) and advocate
staged release with monitoring for misuse.

\subsection{Cache Integrity Defense}\label{sec:cache_defense}

To address the cache injection risk, we present a Cache Integrity
Monitor (CIM) that detects unauthorized KV-cache modifications during
generation. The monitor provides three detection levels with increasing
cost and robustness:

\paragraph{Level 1: Position Fingerprint.}
After encoding, the monitor computes a numerical fingerprint of each
layer's cache tensor (Frobenius norm, mean value, maximum absolute
value, and per-position norm sum). Before each generation step, the
fingerprint of encoding positions is recomputed and compared with
relative tolerance $\epsilon = 10^{-5}$. Any additive perturbation
$\delta$ satisfying $\|\delta\|_F / \|K\|_F > \epsilon$ is detected.
Cost: $O(n)$ per step, where $n$ is the encoding sequence length.

\paragraph{Level 2: Norm Sentinel.}
In addition to Level 1, the monitor tracks per-position L2 norms.
Uniform additive perturbations---the attack demonstrated in
\cref{sec:steering}---change all norms simultaneously, whereas normal
generation only appends one position per step. A change in any
previously-recorded position triggers a violation. Norm distribution
shifts are flagged when the mean exceeds $z > 3.0$ standard deviations
from the prior step's distribution. Cost: $O(n)$ per step.

\paragraph{Level 3: Spectral Sentinel.}
Periodic SVD analysis (every 20 tokens) tracks the trajectory of
Marchenko-Pastur features (signal rank, signal fraction, top SV excess,
spectral gap, norm per token). Anomalous jumps that deviate from the
smooth trajectory expected during normal generation are flagged. This
catches subtle geometric distortions including non-uniform perturbations.
Cost: $O(n^3)$ per check.

\paragraph{Evaluation.}
We evaluated the CIM on two architectures using the same calm direction
vector from the steering experiment: Qwen2.5-7B-Instruct (standard
transformer, 28 layers) and Qwen3.5-27B-Claude-Distilled (hybrid: 16
full-attention + 48 GatedDeltaNet layers). Each was tested under six
scenarios: normal generation (no injection), null injection (zero vector,
confirming the injection code path is a no-op), and calm injection at four
strengths ($\alpha \in \{0.01, 0.1, 1.0, 5.0\}$), with three prompts per
scenario.

\begin{table}[h]
\centering
\small
\begin{tabular}{lccccccc}
\toprule
& & \multicolumn{3}{c}{\textbf{7B (Standard)}} & \multicolumn{3}{c}{\textbf{27B (Hybrid)}} \\
\cmidrule(lr){3-5} \cmidrule(lr){6-8}
\textbf{Scenario} & \textbf{$\alpha$} & \textbf{L1} & \textbf{L2} & \textbf{L3} & \textbf{L1} & \textbf{L2} & \textbf{L3} \\
\midrule
Normal    & 0.00 & 0/3 & 0/3 & 0/3 & 0/3 & 0/3 & 0/3 \\
Null      & 1.00 & 0/3 & 0/3 & 0/3 & 0/3 & 0/3 & 0/3 \\
Calm weak & 0.01 & 3/3 & 3/3 & 3/3 & 3/3 & 3/3 & 3/3 \\
Calm med  & 0.10 & 3/3 & 3/3 & 3/3 & 3/3 & 3/3 & 3/3 \\
Calm std  & 1.00 & 3/3 & 3/3 & 3/3 & 3/3 & 3/3 & 3/3 \\
Calm str  & 5.00 & 3/3 & 3/3 & 3/3 & 3/3 & 3/3 & 3/3 \\
\bottomrule
\end{tabular}
\caption{Cache Integrity Monitor evaluation. FPR: 0/36 = 0\%
(normal + null, both models). TPR: 72/72 = 100\% (all levels, all
strengths, both architectures). With $N = 3$ per cell, Clopper-Pearson
95\% CI for FPR is [0\%, 39\%] per scenario; aggregate 0/36 gives
[0\%, 10\%]. Overhead per 50-token generation: 7B: L1 = 164ms,
L2 = 400ms, L3 = 615ms; 27B: L1 = 207ms, L2 = 522ms, L3 = 1120ms.}
\label{tab:cim_eval}
\end{table}

As shown in \cref{tab:cim_eval}, the monitor detects all tested
uniform additive perturbation injections on both architectures at all
levels, including the weakest strength of $\alpha = 0.01$---two orders
of magnitude below the operational strength used in our steering
experiment. False positive rate is 0/36 across both models (95\% CI:
0--10\%), including the null injection control, which exercises the
full injection code path with a zero vector. True positive rate is
72/72 (95\% CI: 95--100\%). The hybrid architecture poses no challenge: the
monitor correctly processes only the 16 full-attention layers, skipping
GatedDeltaNet layers that lack standard KV-cache tensors.

The per-scenario sample size ($N = 3$) is small. The 0\% FPR point
estimate is encouraging, but the per-scenario Clopper-Pearson 95\% CI
extends to 39\%. A powered evaluation ($N \geq 50$) is needed before
deployment claims. The aggregate 0/36 across all legitimate scenarios
narrows the CI to [0\%, 10\%].

\paragraph{Overhead.}
Level 1 adds approximately 3.3--4.1~ms/token (7B--27B), Level 2 adds
8.0--10.4~ms/token, and Level 3 adds 12.3--22.4~ms/token. For context,
generation speed on 3$\times$RTX 3090 is approximately 100~ms/token (7B)
and 250~ms/token (27B), so Level 1 represents 3--4\% overhead and
Level 3 represents 9--12\%. These costs are non-trivial for latency-sensitive
applications but acceptable for safety-critical deployments where
undetected cache manipulation is the alternative.

\paragraph{Limitations.}
The CIM protects against \emph{post-encoding} injection (the attack
demonstrated here) but does not defend against weight modification,
compromised model code, or injection of adversarial tokens in the input.
Level 1 cannot detect modifications to generation-phase positions
(positions created during generation itself); Level 2 can, via
per-position norm tracking. Only one attack type has been tested (uniform
additive perturbation); manifold-aligned injection that perturbs along
existing singular vectors may evade spectral detection. Cache rollback
requires storing a full copy of the cache, doubling memory during generation.

\paragraph{Release strategy.}
We release the CIM code alongside this paper to ensure the defense is
available concurrently with the attack methodology. The monitor is
implemented as a standalone Python module compatible with any
HuggingFace cache object, requiring no modifications to model code.
